\chapter{Preliminaries, Background, and Related Work}
\label{ch:background}

\section*{Declaration of Contributions}
This chapter was prepared primarily by Utkarsh Kumar and Tanay Kapadia.
Utkarsh Kumar synthesized related research on package ecosystem threats and the MALOSS pipeline,
while Tanay Kapadia contributed dependency confusion background and static source analysis concepts.
Adesh Palkar contributed the generalized threat taxonomy and chapter-level consolidation.

\section{Preliminaries and Basic Notions}

Modern applications depend on external libraries distributed through package registries such as
npm, PyPI, and RubyGems. Let $G=(V,E)$ denote a dependency graph where each node
$v \in V$ is a package and each edge $e \in E$ represents a dependency relation.
A compromise in an upstream node can propagate along outgoing dependency paths to multiple
downstream consumers.

This propagation model explains why software supply chain attacks are ecosystem risks rather than
isolated endpoint events.

\section{Threat Context for Supply Chains}

Traditional cyber defenses are often perimeter-centric. In contrast, supply chain attacks exploit
implicit trust in package origin, maintainer identity, build artifacts, and update channels.
The attack vector is frequently trust itself: a dependency that appears legitimate can execute
malicious logic during install, import, or runtime.

\begin{figure}[htbp]
\centering
\fbox{\parbox{0.86\textwidth}{
\textbf{Conceptual Supply Chain Path:}
Developer publishes package $\rightarrow$ registry distributes package
$\rightarrow$ downstream projects install dependency $\rightarrow$ compromise fans out to users.
An attacker may inject code upstream or exploit naming/version resolution rules.
}}
\caption{Conceptual propagation path in a software supply chain attack.}
\label{fig:supply-chain-path}
\end{figure}

\section{Related Work}

Duan \textit{et al.} presented one of the most influential ecosystem-scale studies on package
manager attacks and proposed MALOSS (MALware detection for Open Source Software), combining
metadata analysis, static analysis, dynamic analysis, and human validation \citep{duan2021ndss}.

Their findings show that malicious packages can persist for long periods and that attack vectors
such as typosquatting and account compromise are recurrent. This supports a data-driven,
continuous-analysis approach rather than one-time manual vetting.

Dependency confusion was later formalized as a practical and severe risk in mixed public-private
dependency environments. Birsan demonstrated real-world exploitation in enterprise contexts
\citep{birsan2021dependencyconfusion}, and Ding \textit{et al.} proposed source tree matching for
systematic risk assessment \citep{ding2024dependency}.

\section{Why Existing Trust Models Fail}

Most ecosystems optimize for package availability and developer convenience.
As a result, they can inherit the following weaknesses:

\begin{itemize}
\item weak provenance guarantees for package origin and release artifacts;
\item delayed detection of malicious behavior hidden in build logic;
\item version-selection policies that can prioritize untrusted public packages.
\end{itemize}

These observations motivate the case study and problem formulation in \cref{ch:casestudy},
followed by in-depth mitigation design in \cref{ch:solutions}.
